\guidechapter{How unwanted sources cancel out}{Read the coefficients, not just the circuit's arrows.}

The source circuit temporarily combines many values. The challenge is to
leave precisely the source meant for each target. The original construction
labels a source and a target with sets of three elements. Their overlap can
have size zero, one, two, or three. Two sets of size three overlap in three
elements only when they are the same set.

For example, $\{1,2,3\}$ and $\{3,4,5\}$ overlap once, while
$\{1,2,3\}$ and $\{2,3,4\}$ overlap twice. The central part of
the complex construction gives coefficient $(q-1)/2$ when the overlap is $q$.
That is $-1/2,0,1/2,1$ across the four possibilities. Side contributions
cancel the two unwanted nonzero cases. The chart shows both contributions
and their sum.

\explfigure{21-intersection-cancellation}{Only the self-contribution survives.}{
The complex coefficients are ordinary signed numbers. The bit construction
instead uses modulo-two arithmetic, where $1+1=0$; its side contribution
cancels the overlap-one case. Both leave net coefficients $0,0,0,1$.}

For a target labeled $S$, every distinct source $T$ therefore has zero
net coefficient, and source $S$ has coefficient one. The target receives
exactly its intended value even though many other values were used internally.
This is stronger than testing one convenient input. It describes every
source's contribution to the output.

\paragraph{Where workspace restoration enters.}
The cancellation table concerns the newly introduced source values. It does
not by itself remove arbitrary values that were already in scratch registers.
The separate restoration schedule cancels those initial values and returns
them. The source proof needs both properties: correct new contribution and
unchanged borrowed data. The register-trajectory figure makes the latter
property visible for one small example.

\takeaway{A crowded network can implement a simple map because unwanted
coefficients cancel. The proof must account for every contribution, including
the values initially occupying its workspace.}
